Привет, меня зовут Кирилл, и в свободное время я вношу вклад в дистрибутив
GNU/Linux под названием Debian.

Моя роль в проекте Debian называется DM (Debian Maintainer)
\url{https://nm.debian.org/person/krekhov/}. Я сопровождаю небольшие и простые
пакеты --- для меня это не только хобби, но и возможность общаться с людьми из
разных стран, культур и сообществ. Среди них есть разработчики, тестировщики,
технические писатели и инженеры, и всех нас объединяет одно --- Debian. Это
прекрасная операционная система, которой я пользуюсь с 2019 года.

Я начал вносить вклад в Debian в 2024 году. Мне очень хотелось оставить свой,
пусть и небольшой, след в истории коммитов интересного мне проекта и принести
пользу. Поначалу мне было трудно разобраться, что к чему: к кому обращаться,
примут ли мои правки, будут ли критиковать за неточности или за что-то ещё. Да,
вначале я нередко спотыкался, но сообщество Debian оказалось довольно
дружелюбным, и я нашёл людей, с которыми мне приятно работать над Debian.

Изначально я начал с документации --- это был проект под названием
\href{https://salsa.debian.org/manpages-l10n-team/manpages-l10n}{manpages-l10n}.
Это проект по переводу man-страниц на разные языки, в том числе на русский. То
есть это переводы справочных страниц команд и программ, которые используются в
Debian и других GNU/Linux-системах. Затем я уже глубже погрузился в пакеты
Debian, разработку, тестирование, сопровождение, сборку, CI/CD и другие
процессы.

Debian --- это не только разработка и техническая работа, но и общение. В
основном оно происходит в почтовых рассылках, Salsa GitLab, BTS (Bug Tracking
System) и IRC-каналах (Internet Relay Chat).

Debian --- это еще и FOSS-культура (Free and Open Source Software), сообщество и
философия, с которыми вы можете ознакомиться здесь:

Social Contract: \url{https://www.debian.org/social_contract}

Philosophy: \url{https://www.debian.org/intro/philosophy}

Code of Conduct: \url{https://www.debian.org/code_of_conduct}

За что я люблю FOSS? За свободу: свободу выбора задач, которые мне интересны.
Я могу углубиться в компиляцию какого-нибудь программного обеспечения, если мне
это интересно именно сегодня, изучать различные флаги GCC (GNU Compiler
Collection), заняться ручным тестированием или написанием тестов, а то и вовсе
открыть GDB (GNU Debugger), начать отлаживать какой-нибудь Core Dump и сообщить
другим инженерам, что мне удалось найти. Меня никто не заставляет делать то,
чего я не хочу: я могу вообще отказаться от дальнейшей работы с багом или чем-то
еще, вежливо написав, что моя работа зашла в тупик и стоит попросить помощи у
кого-то другого (у меня такое редко возникает, но было два случая). Я считаю,
что любой вклад в FOSS важен --- абсолютно любой. Здесь важен каждый.

Я использую последний стабильный выпуск Debian GNU/Linux (Trixie) на amd64.

Мои социальные сети:

- GitHub: \url{https://github.com/krekhovx}

- TelegramChannel: \url{https://t.me/krxnotes}

Данное руководство:

- Репозиторий: \url{https://github.com/krekhovx/debian-way-ru}
